\begin{flushright}
\huge{\textbf{Abstract}}
\end{flushright}

\vspace{0.5cm}

In this project, I implemented and compared several lossless image compression algorithms. I looked at classical entropy methods, spatial and transform techniques, context‑based approaches, clustering, optimization, and even neural networks.

I started with basic entropy coding: Huffman, arithmetic, and LZW. These give you a foundation. Then I tried spatial methods – variable block size compression and run‑length encoding (RLE). The idea is to exploit local repetition in an image. I also implemented a lossy‑plus‑residual scheme: you send a lossy version first, then the difference, so you can reconstruct the original exactly.

Next, I used the JPEG lossless prediction standard with its seven prediction modes. Then I moved to more advanced context‑based algorithms – CALIC and LOCO‑I. They work better because they use gradient information and model the prediction errors.

For transform methods, I implemented the S‑Transform, S+P, and integer wavelet transforms using the lifting scheme. These methods concentrate energy into fewer coefficients, which helps compression.

I also tried clustering for color quantization: K‑means and Gath‑Geva fuzzy clustering. Gath‑Geva can find ellipsoidal clusters of different sizes and densities. Then I used particle swarm optimization (PSO) to search for good color palettes. Finally, I built a fuzzy backpropagation neural network (FBPNN) that combines Gath‑Geva memberships with backpropagation learning.

All the code is in Python. I used NumPy, Matplotlib, Pillow, and GPU‑accelerated frameworks like JAX and PyTorch on an Apple M4. I compared the algorithms with visual outputs, compression ratios, MSE, and PSNR. In the end, the context‑based methods – CALIC and LOCO‑I – gave the best compression, outperforming the traditional ones.

\vspace{1cm}

\noindent\textbf{Keywords:} Lossless Image Compression, Huffman Coding, Arithmetic Coding, LZW, Run-Length Encoding, JPEG Lossless Prediction, CALIC, LOCO-I, S-Transform, S+P Transform, Lifting Wavelet, K-Means Clustering, Gath-Geva Fuzzy Clustering, Particle Swarm Optimization, Fuzzy Backpropagation Neural Network, Context-Based Compression, Transform Domain Compression, GPU Acceleration.